% !PRINTED 143
\weberpage{143}
\begin{center}
\S.\,36.

\gesperrt{Grenzen des Index eines Theilers der symmetrischen
Permutationsgruppe.}
\end{center}

Wir beschliessen diese Betrachtungen mit dem Beweise eines
Satzes über Permutationsgruppen, der durch die Schwierigkeit,
die sein Beweis anfangs bot, eine gewisse Berühmtheit erlangt
hat, und der für die Beurtheilung algebraischer Fragen von
Wichtigkeit ist\,$^{1}$).

Es handelt sich dabei um die symmetrische Permutationsgruppe
 $P$ von $n$ Ziffern und um ihre Theiler von möglichst
kleinem Index. Wir wissen, dass die Gruppe $P$ immer einen
Theiler vom Index 2 hat, nämlich die alternirende Gruppe.
Ausserdem ist noch ein Theiler vom Index $n$ bekannt, der alle
Permutationen von $P$ umfasst, die eine Ziffer ungeändert lassen,
der also intransitiv ist.

Zunächst gilt der folgende Satz:

I. \gesperrt{Der Index eines imprimitiven Theilers von $P$ ist
immer grösser als $n$, und der Index eines intransitiven
 Theilers ist gleich oder grösser
als $n$, und nur dann gleich $n$, wenn der Theiler
eine Ziffer in Ruhe lässt, und die übrigen
$n - 1$ Ziffern auf alle mögliche Arten permutirt.}

Nehmen wir an, es sei $Q$ ein imprimitiver Theiler von $P$
vom Index $j$, und es bestehen $r$ Systeme der Imprimitivität von
je $s$ Ziffern, so dass $n = rs$ ist. Eine Zahl, die der Grad der
Gruppe $Q$ sicher nicht übersteigen kann, erhalten wir, wenn wir
alle Permutationen in jedem einzelnen der $r$ Systeme und dann
noch sämmtliche Permutationen der Systeme abzählen. Der Grad
von $Q$ ist also kleiner oder gleich dem Producte $[\Pi(s)]^r\,\Pi(r)$,
wenn für jede ganze Zahl $n$
\[
\Pi(n) = 1.2.3\dots n
\]
%CONT>>

$^{1}$) \gesperrt{Bertrand}, Journal de Mathématiques, Tome XV (1845). \gesperrt{Serret},
Algèbre supér., Section IV, Chapitre III. C. \gesperrt{Jordan}, Traité des substitutions,
p.\,67. \gesperrt{Netto}, Substitutionentheorie, Capitel VI. Crelle's Journ.,
Bd. 100.
